\section{Evaluation correspondences on expanded configurations}
\label{ev:section}

We construct the extending correspondence required in
\cref{fp:local-algebra} and determine its restriction to each component.
The restriction retains a joint evaluation into a configuration space.
We subsequently prove a product formula which preserves this joint
evaluation when other disconnected target pieces are integrated out.
All virtual classes in this section are ordinary, unreduced classes.

\subsection{The configuration family}

Let $\pi:\mathcal Y\to B$ be a projective semistable degeneration over
a nonsingular pointed curve $(B,0)$, with smooth total space and
\[
  Y_0=X_0\cup_D X_1.
\]
The two sides and $D$ may be disconnected; $D$ is a smooth divisor in
each side and there are no triple intersections.  We fix this coloring
of the sides.  In particular, ``order'' below refers to travel from
$X_0$ to $X_1$ along an expanded interval, and does not impose an
additional ordering on the connected components of $D$.
All constructions can be restricted to an analytic disk about $0$
after they have been made over $B$.

For a finite set $I$ of recorded ordinary markings, write
\[
  \mathcal Q_I=\mathcal Y^{[I]}_\pi,
  \qquad q_I:\mathcal Q_I\longrightarrow\mathcal Y^I_B
\]
for the space of stable expanded configurations.  Its objects are
expansions of $\mathcal Y/B$ with $I$-labelled points in their smooth
locus, such that every exceptional level contains at least one of
these points.  Points are allowed to coincide.  For a smooth pair
$(X,D)$, use similarly $\mathcal Q_I(X,D)=X_D^{[I]}$; its points avoid
the distinguished relative divisor.  We set
$\mathcal Q_\varnothing=B$ and
$\mathcal Q_\varnothing(X,D)=\pt$.

The existence, smoothness, and projectivity of these spaces over the
ordinary products are the configuration-space results
\cite[Propositions 1.2.5 and 1.5.4]{AFConfigurations}.
We record the additional local descriptions that we will use.

\begin{lemma}[Ordered charts and central components]
\label{ev:configuration}
The family $\mathcal Q_I\to B$ is projective and semistable, with
smooth total space.  Its projections to the factors of $\mathcal Y$
are logarithmic maps.  Let $I=I_0\sqcup I_1$ be an assignment of the
recorded marks to the two sides.  The corresponding central component
is canonically
\begin{equation}
\label{ev:component-product}
  Q_{I_0,I_1}
  =\mathcal Q_{I_0}(X_0,D)\times
    \mathcal Q_{I_1}(X_1,D).
\end{equation}
When a side is disconnected, this formula can be restricted to any
prescribed connected target component for each recorded mark.
The aligned configuration space of all marks on a side is retained
in this formula.
\end{lemma}

\begin{proof}
Near a collection of double-locus points, write the fiber product as
\[
  x_i y_i=s,\qquad 1\leq i\leq m=|I|,
\]
omitting smooth coordinates along the double loci.  For a permutation
$\sigma$ of the recorded labels, an ordered chart has coordinates
$a_0,\ldots,a_m$ and equations
\begin{equation}
\label{ev:monomial-chart}
  x_{\sigma(i)}=\prod_{k=0}^{i-1}a_k,
  \qquad
  y_{\sigma(i)}=\prod_{k=i}^{m}a_k,
  \qquad
  s=\prod_{k=0}^{m}a_k.
\end{equation}
These are the charts for the subdivision of the cube by ordered
coordinates.  The simplices are unimodular.  Thus the total space is
smooth and the central fiber is reduced with simple normal crossings.
Factors away from the double locus are treated by making the
appropriate coordinates invertible.  The projections are monomial
in these charts and hence logarithmic.

The charts have an intrinsic interpretation: put the recorded points
on a common expansion and contract precisely those exceptional levels
containing none of them.  Ratios of normal coordinates on a retained
level record the relative positions of its points.  Changing a local
equation of $D$ multiplies such a coordinate by a unit and changes the
normal coordinate by the transition function of its divisor line
bundle.  Consequently the descriptions glue globally.  They are the
configuration spaces above, whose projectivity is supplied by the
cited construction.

The divisor $a_j=0$ in \eqref{ev:monomial-chart} puts the first $j$
marks on one side of a cut and the remaining marks on the other.
Gluing the distinguished divisors of the two expanded pairs on the
right of \eqref{ev:component-product} gives a configuration in this
component.  All exceptional levels remain occupied, so this gluing
does not require a further contraction.  Conversely, a stable
configuration in the indicated component has a unique cut with
exactly $I_0$ on the $X_0$ side.  Two distinct such cuts would enclose
an exceptional level containing no recorded point.  Cutting there
recovers the two factors of \eqref{ev:component-product}.  Scheme
theoretically, on $a_j=0$ the remaining coordinates separate into
$a_0,\ldots,a_{j-1}$ and $a_m,\ldots,a_{j+1}$.  These are the
ordered charts of the left relative configuration and the right
relative configuration with reversed order.  The smooth coordinates
along $D$ accompany their respective labels.  Further vanishing
coordinates describe deeper strata inside the two factors.  This
also proves that the constructions commute with base change and
are inverse as morphisms.  There is no additional fiber product over
$D$: the divisors of the targets are glued, whereas the recorded
interior positions have independent projections to $D$.
If $I_0$ or $I_1$ is empty, its relative configuration is the
unexpanded pair with no recorded positions, represented by $\pt$.

For disconnected sides the collapsed evaluation of a point determines
its connected component of $X_i$.  Restricting this locally constant
choice gives the last assertion.  In particular, no independent
expansion is chosen for each connected component within a side.
\end{proof}

Forgetting any subset of recorded positions defines a morphism between
the corresponding configuration spaces: after forgetting, contract
the levels that have become empty.  This is the target stabilization
used throughout this section.  It does not forget or stabilize any
marking of the source curve of a stable map.

\subsection{The extending virtual correspondence}

Fix discrete map data and a finite set of ordinary markings containing
$I$.  The source may be disconnected.  Let $\mathcal M$ be the
fixed-data twisted transverse-map stack of Abramovich--Fantechi over
the universal expansion stack; its coarse maps are predeformable.
Throughout the correspondence, cap, and sewing arguments we use this
expansion convention.  Its levels
are globally aligned even when $D$ is disconnected, exactly as in
the configuration spaces of \cref{ev:configuration}; we do not
identify this stack with a presentation allowing independent
expansion lengths on the connected components of $D$.  In this paper
``ordinary relative theory'' means scheme-target invariants with
ordinary coarse-source insertions computed in this aligned formalism.
Properness of the coarse predeformable-map stack in this universal
model follows from \cite[Corollary~2.2.9]{AFConfigurations}, and
twisted-map properness and relative virtual classes follow from
\cite[Theorem~3.26 and Section~4.7]{AF2016}.  We retain stable
unmarked rootless components in this labelled theory; the
common-expansion obstruction comparison below applies to them
directly.

The descendant class is the cotangent line of the coarse source at
an ordinary marking on the expanded map.  Stabilization over the
unexpanded target, retaining all ordinary and relative markings, is
an isomorphism near each ordinary marking
\cite[Definition~4.1 and Section~4.3]{AF2016}.  Thus this line is
also the pullback of the ordinary stable-map cotangent line used
in the comparison of relative theories
\cite[Corollary~1.1.3 and Section~2.3]{AMWComparison}.
Target-level stabilization, cutting, and the group separation below
preserve the same neighborhood and hence this cotangent line;
see also \cite[Section~5.10]{AF2016} for gluing.  On a smooth fiber
there are no target expansions, so this is the ordinary map
cotangent.  When smooth
curve data have several specializations, take their full finite sum
with the specified extending numerical degree weights.  No individual
lifted homology class is chosen across different smoothings.  The
usual bounded-degree conditions make these constructions finite
coefficient by coefficient.

An ordinary marking lies in the smooth interior of the expanded
target.  Retain the images of the markings in $I$ and stabilize these
target positions.  This gives a morphism
\begin{equation}
\label{ev:lift}
  e_I:\mathcal M\longrightarrow\mathcal Q_I.
\end{equation}
For twisted expansions we first take their coarse expansion
coordinates.  Ordinary interior positions have no ambiguity under
this operation.  The morphism is proper: $\mathcal M$ is proper over
$B$ and $\mathcal Q_I$ is separated over $B$.  Target stabilization
is defined for families, either by the preceding configuration
construction or by multiplying the smoothing parameters at the
contracted levels in \eqref{ev:monomial-chart}.

Let
\[
  \Theta=\prod_{j}\psi_j^{b_j}\prod_{j\notin I}\ev_j^*\gamma_j
\]
denote the ordinary descendant powers and preintegrated insertions.
The $\psi_j$ are cotangent classes on maps.  Initially suppose that
the $\gamma_j$ are extending algebraic classes.  Push forward the
virtual class over the full base:
\begin{equation}
\label{ev:total-class}
  \xi_I=(e_I)_*\bigl(\Theta\cap[\mathcal M]^{\vir}\bigr)
  \in A_*(\mathcal Q_I)_\Q,
  \qquad
  \Gamma_I=\operatorname{cl}(\xi_I).
\end{equation}
We use the identification of Chow homology with Chow cohomology on
the smooth total space to write its cycle class as $\Gamma_I$.
If the correspondence has codimension $c$, then $\Gamma_I$ has
cohomological degree $2c$ and Tate type $(c,c)$.
Restriction to a noncentral fiber of \eqref{ev:total-class} is the
usual descendant evaluation correspondence.  This follows from
virtual base change and proper pushforward, since over that fiber
there are no exceptional target levels.

The construction is made before any nonalgebraic constant-source
inputs are contracted.  More explicitly, expose such inputs as
cohomology slots of their fixed smooth projective sources and use
the algebraic correspondences extending over the total family.
The virtual
pushforward with these slots exposed is a morphism of Hodge
structures with its prescribed Tate shift.  Contracting with the
fixed inputs afterward gives the required preintegrated
correspondence.  If desired, one can record their target positions
as additional factors, apply the same construction, and push forward
while contracting the source slots.  The projection formula makes
the two procedures equal.  The number of positions used in the later
local calculation is therefore the number in $I$, not the number of
ordinary preintegrated markings.  General cohomology inputs in these
fixed slots follow by linearity.

\begin{lemma}[Labelled cut in the aligned expansion]
\label{ev:aligned-cut}
Fix bounded discrete data for labelled connected source components
in $\mathcal M$, permitting their disjoint union, and choose one cut
of the globally aligned target expansion.  Let $\eta$ range over the
resulting splittings with ordered labels on their $r$ paired relative
roots of contact
orders $c_1,\ldots,c_r$.  The side source graphs may contain rootless
vertices or be empty.  Let $\mathcal M_{\eta,i}$ denote the relative
map stack of the \emph{whole} side source graph over one shared side
expansion.  On the normalized target cut, the virtual boundary class
is the sum, over these labelled splittings, of the gluing pushforwards
of
\begin{equation}
\label{ev:aligned-cut-formula}
 \frac{\prod_{j=1}^r c_j}{r!}
 \Delta_{D^r}^{!}
 \bigl([\mathcal M_{\eta,0}]^{\vir}
       \times[\mathcal M_{\eta,1}]^{\vir}\bigr).
\end{equation}
There is no additional factor for a source component or a connected
component of $D$.  For $r=0$, the product and diagonal are the identity.
The formula is an equality of virtual classes before evaluation
pushforward.  The $r!$ forgets the ordered root labels; the finite
quotient by permutations of identical source components is taken on
the whole indexed sum, not again on each root-labelled side factor.
\end{lemma}

\begin{proof}
The untwisted target stack with a chosen global cut normalizes its
normal-crossings boundary with pure degree one
\cite[Proposition~7.4.1]{ACFW}.  Base-changing the map stack and its
relative obstruction theory to this normalization gives the virtual
cut class by the virtual pushforward comparison used in
\cite[Sections~5.4--5.6]{AF2016}.  This target-stack statement does
not depend on the source graph or on the number of components of $D$.
For a twist index $R$, the reduced twisted split stack has degree
$1/R$ over the corresponding nonreduced boundary; hence this step
contributes $R$ to the virtual class.

Normalize the source at its nodes mapping to the selected cut.  A
rootless component remains wholly on one side, including when it
maps into an expanded level.  On fixed labelled data the split-map
stack is the fiber product of the two relative-map stacks over the
matching root evaluations and the \emph{single} target-gluing stack.
Relative to this common target-expansion base, the AF map-deformation
complex on a disjoint union is a direct sum.  Its source-normalization
triangle has exactly the normal complex of the matching-root diagonal
as the difference between split and glued complexes
\cite[Proposition~5.16]{AF2016}; there is no term at a rootless
component.  Thus the compatible obstruction theories give the
refined diagonal pullback in \eqref{ev:aligned-cut-formula}.  The
common-refinement comparison for disconnected source maps
\cite[Proposition~4.14 and its proof]{AF2016} keeps one expansion per
side; it does not duplicate the target deformation for each source
component.  Although the published degeneration formula indexes
connected glued graphs, its target normalization and this labelled
source-normalization comparison do not use that condition.  When
both sides are nonempty, connectedness ensures that every side
vertex has a root.  Here rootless vertices and empty sides are
retained, and source-component labels are kept until their finite
quotient is taken.

The universal split target over the two side expansion stacks is one
$\mu_R$-gerbe, of degree $1/R$
\cite[Proposition~7.4.2]{ACFW}; this cancels the preceding $R$
\cite[Lemma~5.15]{AF2016}.  The separate comparison from rigidified
root evaluations to the ordinary diagonal has degree
$\prod_jc_j$ \cite[Proposition~5.18]{AF2016}.  These are rootwise
factors even if the roots lie on different connected components of
$D$; the target-twist factor remains global.

If a side source graph is empty, every positive-length expansion on
that side has an unstabilized level-scaling automorphism.  Its stable
relative-map stack is therefore the point represented by the
unexpanded pair, with zero relative obstruction complex and virtual
class $[\pt]$.  For an empty root profile choose $R=1$; the target
gerbe is trivial, $D^0=\pt$, and both contact and root-permutation
factors are one.  A nonempty rootless source component on a bubble
is retained on its side and is not confused with this empty unit.
Finally, forgetting the ordered root labels is an $r!$-cover and
gives the denominator in \eqref{ev:aligned-cut-formula}.  Quotienting
the fixed source-component labels by permutations of identical
components gives their separate stack weights.
\end{proof}

\begin{proposition}[Restriction with recorded assignments]
\label{ev:restriction}
The class $\Gamma_I$ is an actual extending correspondence on the
smooth total configuration family.  Let
$i_\epsilon:Q_\epsilon\hookrightarrow\mathcal Q_I$ be a central
component, with the recorded side and connected-component assignment
$\epsilon$.  Its restriction $i_\epsilon^*\Gamma_I$ is given by the
virtual degeneration formula restricted to splitting data with that
assignment, retaining on each side the joint relative evaluation
into the corresponding factor of \eqref{ev:component-product}.
This is an equality of correspondence classes on $Q_\epsilon$
before further cohomology insertions are applied.

Concretely, in the scheme case let $\eta$ range over admissible
splittings with that assignment, with $r$ labelled matching relative
roots of orders $c_1,\ldots,c_r$.  Let
$\mathcal M_{\eta,0}$ and $\mathcal M_{\eta,1}$ be the relative
moduli, and let
\[
  G_\eta=\mathcal M_{\eta,0}\times_{D^r}\mathcal M_{\eta,1}.
\]
The root matching is the refined diagonal pullback.  With the usual
labelled-root convention, the restriction is the cycle class of
\begin{equation}
\label{ev:restriction-formula}
 \sum_{\eta\text{ of assignment }\epsilon}
 \frac{\prod_{j=1}^r c_j}{r!}\,
 (e_\eta)_*
 \left(
   \Theta_\eta\cap
   \Delta_{D^r}^{!}
   \bigl([\mathcal M_{\eta,0}]^{\vir}
          \times[\mathcal M_{\eta,1}]^{\vir}\bigr)
 \right).
\end{equation}
Here $e_\eta:G_\eta\to Q_\epsilon$ retains the two aligned relative
evaluations; it is allowed to have image on the boundary of
$Q_\epsilon$.  Source-component labeling or quotient conventions are
held fixed throughout.  Equivalently, summing over unlabelled root
data replaces $r!$ by the usual root-permutation stabilizer.
For $I=\varnothing$ the assertion is the ordinary special-fiber
virtual degeneration formula.
\end{proposition}

\begin{proof}
We have already constructed the extending class.  We prove the
assignment-specific statement by a Cartier divisor calculation,
then apply virtual splitting on each resulting boundary divisor.

On a chart of the stack of expanded targets let
$t_0,\ldots,t_\ell$ be the untwisted smoothing parameters, so that
$s=t_0\cdots t_\ell$.  At an ordinary recorded point the normal
coordinate along its level is a unit.  Order the recorded points by
their levels.  Evaluation into \eqref{ev:monomial-chart} has the
following form: each $e_I^*a_j$ is a unit times the product of the
$t_k$ belonging to the interval between the two successive recorded
levels, with the two end intervals used for $j=0,m$.  Points on the
same level give unit ratios.  Consequently, as divisor line bundles
with their defining sections,
\begin{equation}
\label{ev:cut-divisor}
  e_I^*(a_j=0)
  =\sum_{k\in J_j}(t_k=0),
\end{equation}
where $J_j$ consists exactly of the cuts putting the recorded marks
on the sides specified by $a_j=0$.  There is order one on each of
these divisors in untwisted target coordinates and order zero on the
other cut divisors.  Empty unrecorded levels simply contribute more
factors to the same interval product.  On a normalized splitting
divisor, the connected component of a side containing a recorded
mark is also discrete.  Selecting that choice gives precisely the
refinement from the side assignment to $\epsilon$.

Here the notation in \eqref{ev:cut-divisor} means the factorization
of the pulled-back section and its line bundle.  It does not assert
that this section is regular on every component of $\mathcal M$.
The virtual boundary intersections below are refined pullbacks of
the boundary divisors on the expansion stack.  Their additivity
under this factorization therefore also applies to components of
the map space supported over $s=0$.

The equations change by units under normal-coordinate transitions,
so \eqref{ev:cut-divisor} is an equality of the corresponding divisor
line bundles and sections, not just an equality on an open set of
maps.  In particular, it controls refined intersection also when the
evaluation image lies in a deeper central stratum.  The map from
each normalized cut divisor to $Q_\epsilon$ is obtained by cutting
the expansion and stabilizing the recorded positions on its two
sides.  The uniqueness argument in \cref{ev:configuration} shows
that this is the map occurring in
\eqref{ev:component-product}.

Refined Cartier pullback commutes with proper pushforward.  Applied
to \eqref{ev:total-class}, it identifies the homological class
corresponding to $i_\epsilon^*\Gamma_I$ with the pushforward of the
virtual intersection with the left side of
\eqref{ev:cut-divisor}.  Intersect the right side one boundary at a
time.  The virtual splitting identities in the proof of the
degeneration formula identify each such intersection with the
matching relative virtual classes and their refined root diagonal.
The required identity for disconnected sources and an empty side is
\cref{ev:aligned-cut}, applied before evaluation pushforward.  The
additional assertion beyond the numerical degeneration formulas
\cite[Sections~5.3--5.9]{AF2016} and \cite[Theorem~2.3]{ABPZ} is the
retained correspondence with its recorded assignment and the full
labelled source graph, including rootless components.

A target level is retained when stabilized by the union of its
source components; after cutting, each side is stabilized separately.
Bounded expanded-map stacks supply properness and finite type.  The
target-cut normalization and Cartier equation
\eqref{ev:cut-divisor} depend only on the target expansion and the
recorded assignment, so they are unchanged by the source graph.

Passing to unlabelled disconnected sources divides both sides by the
finite permutation group of isomorphic components; intrinsic map
automorphisms remain in the stack virtual classes.  Two identical
closed components on one side, for example, contribute $1/2!$,
while the two labelled assignments with one on each side cancel
that quotient.  Labelling the relative roots and dividing by $r!$
gives the convention of \eqref{ev:restriction-formula}.  This proves
the formula as a correspondence class before further cohomology
insertions or numerical integration, with the aligned evaluation in
$Q_\epsilon$ retained.

For clarity about multiplicities, a twisted node parameter $\tau$
has untwisted parameter $t=\tau^r$.  Pulling back the untwisted
component equation therefore gives its twisting multiplicity.
The gerbe degree in gluing and the root-pairing normalization in the
twisted splitting identity combine with it to give the ordinary
contact coefficient in \eqref{ev:restriction-formula}.  One must
not replace this calculation by an assertion of multiplicity one
in source-node or twisted coordinates.  We have used multiplicity
one only in the untwisted target boundary equation
\eqref{ev:cut-divisor}.

Cutting and gluing a target leave the curve unchanged in a
neighborhood of every ordinary marking.  Hence its ordinary
cotangent line and all powers of that line restrict as used in
$\Theta_\eta$.  Preintegrated extending classes restrict by the
projection formula.  For constant-source Hodge inputs the equality
is first made with the source slots exposed and then contracted,
as described above.  Finally, with no recorded positions the sole
component equation is $s=0$ and all cuts occur; the same proof gives
the usual formula.
\end{proof}

\begin{remark}
\label{ev:refinement-not-numerical}
The enhancement in \cref{ev:restriction} is the calculation
\eqref{ev:cut-divisor} and its use before evaluation pushforward.
The standard virtual degeneration theorem supplies the splitting
of each boundary class.  A numerical formula alone would not
determine the restriction to an individual component of the
resolved configuration family.
\end{remark}

\begin{example}[A diagonal near the double locus]
\label{ev:diagonal-example}
For two recorded points a chart is
\[
 x_1=a,\quad x_2=av,\quad y_1=vb,\quad y_2=b,\quad s=avb.
\]
The lifted diagonal of equal points is $v=1$.  It misses the
mixed-assignment component $v=0$, and the other ordered chart gives
the same statement for the opposite mixed assignment.  Thus the
degree-zero constant-map diagonal has zero restriction on those
components.  For three equal recorded points all intermediate
ratios in \eqref{ev:monomial-chart} equal one.  This illustrates
why restrictions must be computed on the configuration resolution,
and is consistent with \cref{ev:restriction} for a stable constant
three-marked genus-zero map.
\end{example}

\subsection{Separating groups of relative maps}

We next consider a fixed pair $(X,D)$, with a parametrized original
target $X$.  This is relative theory, not rubber theory.  Fix
discrete relative data $\Xi$ for a disconnected source and divide
its connected components into groups
\[
  \Xi=\Xi^{(1)}\sqcup\cdots\sqcup\Xi^{(u)}.
\]
A group may itself be disconnected.  The groups and all source,
ordinary-marking, and relative-root labels are fixed in this
statement.  Write $\mathcal K_\Xi$ for maps on a common expansion
and $\mathcal K_{\Xi^{(a)}}$ for the separately stabilized group.
There is a separation morphism
\begin{equation}
\label{ev:separation}
 \varepsilon:\mathcal K_\Xi\longrightarrow
 \prod_{a=1}^u\mathcal K_{\Xi^{(a)}}.
\end{equation}
For each group choose a subset $I_a$ of its ordinary marks and let
$e_a$ be their joint evaluation into
$\mathcal Q_{I_a}(X,D)$.  An empty $I_a$ gives the constant map to
$\pt$.

\begin{proposition}[Group product with retained evaluations]
\label{ev:groups}
The separation morphism satisfies
\begin{equation}
\label{ev:group-vfc}
  \varepsilon_*[\mathcal K_\Xi]^{\vir}
  =\mathop{\times}_{a=1}^u[\mathcal K_{\Xi^{(a)}}]^{\vir}.
\end{equation}
Ordinary cotangent classes and joint evaluations of the retained
positions are preserved under separation.  Thus if $\Theta_a$
is a product of ordinary descendants and insertions on group $a$,
and $e=(e_1,\ldots,e_u)\circ\varepsilon$, then
\begin{equation}
\label{ev:group-correspondence}
 e_*\left(\prod_a\Theta_a\cap[\mathcal K_\Xi]^{\vir}\right)
 =\mathop{\times}_{a=1}^u
   (e_a)_*\left(\Theta_a\cap
           [\mathcal K_{\Xi^{(a)}}]^{\vir}\right).
\end{equation}
The same assertion holds after grouping by connected components
of a disconnected target and retaining the aligned evaluation of
any chosen group.  Factors with no recorded positions can therefore
be preintegrated as ordinary relative invariants.  The identities
with algebraic insertions hold in Chow groups; for arbitrary
cohomology insertions we apply the cycle class map and the
cohomological projection formula.
\end{proposition}

\begin{proof}
For groups consisting of single connected source components,
\eqref{ev:group-vfc} is the relative disconnected product rule
\cite[Proposition 4.14]{AF2016}.  We describe its comparison with
groups retained, since the additional evaluation assertion cannot
be obtained merely by quoting numerical factorization.

Let $\mathcal T$ denote the stack of expansions of the pair, with
the twisting choices used in the transverse formulation when
needed.  Form the auxiliary stack $\mathcal T'$ of one common
expansion $\mathcal X$ and partial contractions
\[
   \mathcal X\longrightarrow\mathcal X_a,
   \qquad 1\leq a\leq u,
\]
where each $\mathcal X_a$ is an expansion, and every exceptional
level of $\mathcal X$ is retained by at least one contraction.
The twists and partial untwistings are included as in
\cite[Sections~4.7.3--4.7.4]{AF2016}; alternatively they may be put under
a common dominating twisting choice.  The comparison with the
product expansion stack has pure degree one: both stacks have the
same dense locus of unexpanded targets.  The forgetful map from
$\mathcal T'$ to the common expansion stack is the same local
etale comparison as in that construction.

The square
\begin{equation}
\label{ev:group-square}
 \begin{matrix}
  \mathcal K_\Xi &\xrightarrow{\ \varepsilon\ }&
        \prod_a\mathcal K_{\Xi^{(a)}}\\
  \big\downarrow &&\big\downarrow\\
  \mathcal T'&\longrightarrow&\mathcal T^{u}
 \end{matrix}
\end{equation}
is cartesian in the transverse expanded-map comparison.  Indeed,
given maps from the separate groups and a specified common
refinement, pull each map back along
$\mathcal X\to\mathcal X_a$.  Transversality makes this pullback a
curve with the canonical trivial cylinders at the levels contracted
in $\mathcal X_a$.  The markings lift uniquely, and the disjoint
union of the lifted groups is a map to the common expansion.
The requirement that each level is retained by at least one group
gives stability of the combined map.  Conversely, separating a map
and contracting the levels unstable for an individual group recovers
this construction.  No step uses connectedness inside a group.

The relative map obstruction complex on a disjoint union is the
direct sum of the complexes on its groups.  Under the partial curve
contraction $c:C'_a\to C_a$ above, the contracted components are
trivial rational cylinders and
$Rc_*\mathcal O_{C'_a}=\mathcal O_{C_a}$.  The target-relative
complex in the transverse comparison pulls back to the
corresponding complex on $C'_a$.  Projection formula therefore
identifies the relative obstruction theory upstairs with the
pullback of the direct-sum theory on the product.  This is exactly
the obstruction-theory comparison of
\cite[Lemmas 4.16--4.18]{AF2016}, with a disjoint union in place
of each individual connected source.  Its degree-one virtual
pushforward argument applied to \eqref{ev:group-square} proves
\eqref{ev:group-vfc}.

Consider an ordinary marking belonging to group $a$.  A level
containing its image cannot be trivial for that group: the
nontrivial scaling of a fiber of that level cannot fix an interior
marked point while preserving the map.  Consequently the curve
contractions used in separation are isomorphisms near all ordinary
marks of the retained group.  They preserve their cotangent lines.
They also preserve their evaluation positions on the remaining
levels.  Stabilizing these recorded positions before or after
separation gives the same configuration, because in either case
the final expansion contracts precisely the levels empty of those
positions.  Thus $e$ and each $\Theta_a$ factor through
\eqref{ev:separation}.  Projection formula and proper pushforward
give \eqref{ev:group-correspondence}.

The source exposition of the product comparison assumes every
connected component carries a marking or a relative root, expressly
for simplicity \cite[Section~4.7.1]{AF2016}.  Its proof of the
product-class identity extends directly to stable labelled components
with neither.  A level occupied by a positive-degree map is retained
when contracting it would destroy transversality or stability.  A
rootless closed component of collapsed degree zero cannot have a
nonconstant fiber-only part: a
complete nonconstant curve in a ruled fiber meets its boundary, and
predeformability propagates this to a root or positive collapsed
degree.  Thus the remaining unmarked component is constant and has
genus at least two by stability.  The cartesian common-refinement
square, etale expansion comparison, and direct-sum obstruction theory
in \cite[Lemmas~4.16--4.18]{AF2016} depend on these stable maps and
their universal curves, not on the presence of a mark or root.
Their degree-one virtual pushforward proves
\eqref{ev:group-vfc} for these components as well.  The empty
source graph is the unit $[\pt]$.

Finally, a connected source curve mapping into a disconnected
target lies in one connected target component.  We may collect all
such source components into the desired target groups.  The same
argument applies, keeping one expansion and one aligned evaluation
for each group until \eqref{ev:group-correspondence} is used.
With cohomology inputs of odd parity, the products and permutations
in this argument are taken in the graded tensor category.  This
inserts the usual Koszul signs and does not change any of the
virtual pushforward equalities.
\end{proof}

\subsection{The correspondence data used in sewing}

The component formula of \cref{ev:restriction} retains the aligned
relative evaluations required by the local first-page operations.
If a disconnected target piece contains none of their recorded
positions, \cref{ev:groups} permits its ordinary relative data to
be integrated out while preserving the retained correspondence.
The statement remains valid when that piece carries arbitrarily
many ordinary preintegrated auxiliary labels, since those labels
occur only in its factor of \eqref{ev:group-correspondence}.

Both assertions hold for fixed discrete data and hence for any
finite coefficientwise linear combination allowed by the degree
and genus bounds.  They impose no extra test of connectedness in
a smoothed target.  Each relative root matching remains the
ordinary diagonal contraction, with its contact and automorphism
weights; the smooth-theory tensor contractions of
\cref{fp:local-algebra} are performed on the retained correspondence
after this virtual splitting.  This is the form needed at an idle
switch in \cref{sew:transport}.
